\section*{第 \ref{ch:b3:galois} 章——域扩张与
伽罗瓦理论}

\begin{solution}{exo:b3:galois:1}
$\sqrt3 \notin \Q(\sqrt2)$：若 $\sqrt3 = a + b\sqrt2$（$a, b
\in \Q$），平方得 $3 = a^2 + 2b^2 + 2ab\sqrt2$，故 $ab =
0$；$b = 0$ 时 $\sqrt3$ 为有理数，$a = 0$ 时得 $\sqrt6 = 2b
\in \Q$——两者均不成立（标准的素因数分解论证）。
因此 $[\Q(\sqrt2,\sqrt3) : \Q(\sqrt2)] = 2$，由塔定理
得次数为 $4$。

设 $\gamma = \sqrt2 + \sqrt3$，则 $\gamma^2 = 5 + 2\sqrt6$，且
$(\gamma^2 - 5)^2 = 24$，故 $\gamma$ 以 $X^4 - 10X^2 + 1$ 为零化多项式。
又 $\gamma^3 = 11\sqrt2 + 9\sqrt3$，故 $\gamma^3 - 9\gamma
= 2\sqrt2$，于是 $\sqrt2 \in \Q(\gamma)$，进而 $\sqrt3 = \gamma -
\sqrt2 \in \Q(\gamma)$，故 $\Q(\gamma) = \Q(\sqrt2,\sqrt3)$，其次数为
$4$。该零化四次多项式的次数与
\omterm{def:b3:galois:algebraic}{极小多项式}的次数相同，故它\emph{就是}\omterm{def:b3:galois:algebraic}{极小多项式}：$X^4 -
10X^2 + 1$（特别地，它在 $\Q$ 上不可约）。
\end{solution}

\begin{solution}{exo:b3:galois:2}
$X^3 - 2$ 在 $\C$ 中的三个根为 $\alpha, j\alpha,
j^2\alpha$，其中 $j = \eu^{2\iu\pi/3}$。映射 $\alpha \mapsto
j\alpha$ 给出一个 $\Q$-嵌入 $\Q(\alpha) \to \C$
（见\cref{thm:b3:galois:simple}：二者都生成次数为 $3$ 的扩张，
且\omterm{def:b3:galois:algebraic}{极小多项式}相同），其像 $\Q(j\alpha)
\not\subseteq \R$ 不同于 $\Q(\alpha) \subseteq \R$：
故 $\Q(\alpha)/\Q$ 不是正规扩张。\omterm{thm:b3:galois:splitting}{分裂域}为 $L =
\Q(\alpha, j)$，且 $[L:\Q] = [L:\Q(\alpha)]\,[\Q(\alpha):\Q] =
2 \cdot 3 = 6$（$j$ 满足 $X^2 + X + 1$，它在
实域 $\Q(\alpha)$ 上不可约）。$\Q(\alpha)$ 的任一自同构
必将 $\alpha$ 映为 $X^3 - 2$ 的根，且该根\emph{位于}
$\Q(\alpha) \subseteq \R$ 之中：只有 $\alpha$ 符合，于是
$\operatorname{Aut}_\Q(\Q(\alpha)) = \{\mathrm{id}\}$，阶为
$1 < 3$.
\end{solution}

\begin{solution}{exo:b3:galois:3}
$X^2 + 1$ 在 $\mathbb F_3$ 中无根（$0, 1, 2 \mapsto 1, 2,
2$），故 $\mathbb F_9 = \mathbb F_3[X]/(X^2+1)$ 为域，
含 $9$ 个元素；记 $\omega = \bar X$，$\omega^2 = -1$。
$\mathbb F_9^\times$ 是阶为 $8$ 的循环群；$\omega$ 的阶为
$4$，而 $1 + \omega$ 可作生成元：$(1+\omega)^2 = 1 + 2\omega +
\omega^2 = 2\omega$, $(1+\omega)^4 = 4\omega^2 = \omega^2 = -1
\ne 1$：阶为 $8$。

在 $\mathbb F_2$ 上——次数为 $1$ 的：$X$，$X + 1$；次数为 $2$ 的：
$X^2 + X + 1$（其余三个二次式都有根）；次数为
$3$：$X^3 + X + 1$ 与 $X^3 + X^2 + 1$（在 $\mathbb
F_2$ 中无根；其余六个三次式都有根）。验证：
\[
(X^3{+}X{+}1)(X^3{+}X^2{+}1) = X^6 + X^5 + X^4 + X^3 + X^2 + X
+ 1,
\]
又 $X(X{+}1)(X^6 + \dots + 1) = X(X^7 + 1) = X^8 + X = X^8 -
X$ 在 $\mathbb F_2$ 上——恰为次数整除
$3$ 的全体不可约多项式之积，与\cref{exo:b3:galois:6}的预言一致（次数 $2$
不出现：$2 \nmid 3$）。
\end{solution}

\begin{solution}{exo:b3:galois:4}
(a) 模 $7$：$3$ 的幂依次为 $3, 2, 6, 4, 5, 1$，阶为
$6$，为生成元；原根为 $3^k$，其中
$\gcd(k, 6) = 1$，即 $3$ 与 $3^5 = 5$。模 $11$：$2$ 的幂为
$2, 4, 8, 5, 10, 9, 7, 3, 6, 1$，为生成元；原根为
$2^k$，$\gcd(k, 10) = 1$：$2, 2^3 = 8, 2^7 = 7, 2^9 = 6$。

(b) 记 $x = g^k$，其中 $g$ 为生成元
（见\cref{thm:b3:galois:cyclic}）。于是 $x$ 是平方元当且仅当 $k$
为偶数（平方元即 $g^{2l}$，且 $g^{2l} = g^{k}$ 当且仅当 $k
\equiv 2l \bmod p-1$，此方程有解当且仅当 $k$ 为偶数，$p - 1$ 为
偶数）。又 $x^{(p-1)/2} = g^{k(p-1)/2} = 1$ 当且仅当 $(p-1) \mid
k\frac{p-1}2$ 当且仅当 $k$ 为偶数：两个条件一致。对 $x =
-1 = g^{(p-1)/2}$：它为平方元当且仅当 $\frac{p-1}2$ 为偶数，即
$p \equiv 1 \pmod 4$——再次见\cref{pb:b3:rings:1}。
\end{solution}

\begin{solution}{exo:b3:galois:5}
在 $\bar{\mathbb F}_p$ 中，$\mathbb F_{p^k} = \{x : x^{p^k} =
x\}$ 为 $F^k$ 的不动点集。交集 $\mathbb F_{p^m}
\cap \mathbb F_{p^n}$ 被 $F^m$ 与 $F^n$ 固定，从而被
$F^{\gcd(m,n)}$ 固定（$\gcd = am + bn$：对不动元，$F^{am +
bn} = (F^m)^a(F^n)^b$ 作用平凡——指数可取为
正数，由周期性）；故其含于 $\mathbb
F_{p^{\gcd(m,n)}}$，反之它亦含于两者
（见\cref{thm:b3:galois:finitefields}(3)）。复合域 $\mathbb
F_{p^m}\mathbb F_{p^n}$：任一同时包含二者的域，其次数被
$m$ 与 $n$ 整除，从而被 $\operatorname{lcm}(m,n)$ 整除；而
$\mathbb F_{p^{\operatorname{lcm}}}$ 同时包含二者，它即为
该复合域。若 $m \mid n$，则 $\operatorname{Gal}(\mathbb
F_{p^n}/\mathbb F_{p^m})$ 的元素为 $F$ 的幂中
使 $\mathbb F_{p^m}$ 不变者，即 $F^m$ 的幂：循环群，阶为 $n/m$，
由 $F^m \colon x \mapsto x^{p^m}$ 生成（阶的计算见
\cref{thm:b3:galois:finitefields}(2)).
\end{solution}

\begin{solution}{exo:b3:galois:6}
$X^{q^n} - X$ 可分（导数为 $-1$），根集为
$\mathbb F_{q^n}$。设 $P$ 为次数为 $d$ 的首一不可约多项式。
若 $d \mid n$，则 $\mathbb F_q[X]/(P) \cong \mathbb F_{q^d}
\subseteq \mathbb F_{q^n}$，故 $P$ 有根 $\alpha \in
\mathbb F_{q^n}$；$\alpha^{q^n} = \alpha$，且 $P =
\pi_\alpha$ 整除 $X^{q^n} - X$。若 $P \mid X^{q^n} - X$，则
其一根 $\alpha \in \mathbb F_{q^n}$ 生成 $\mathbb F_{q^d}
\subseteq \mathbb F_{q^n}$，故 $d \mid n$
（见\cref{thm:b3:galois:finitefields}(3)）。不同的不可约多项式
互素，且其乘积可分：每个 $P$ 出现的指数恰为
$1$，且 $X^{q^n}-X$ 的每个根都是
其\omterm{def:b3:galois:algebraic}{极小多项式}的根，因式分解成立。比较
次数：$q^n = \sum_{d\mid n} d\,I_d(q)$。

因此 $I_1 = q$；$q^2 = I_1 + 2I_2$ 给出 $I_2 =
\frac{q^2 - q}2$；$q^3 = I_1 + 3I_3$ 给出 $I_3 = \frac{q^3 -
q}3$；$q^4 = I_1 + 2I_2 + 4I_4$ 给出 $I_4 = \frac{q^4 -
q^2}4$。存在性：$nI_n = q^n - \sum_{d \mid n,\, d < n}
dI_d \geq q^n - \sum_{d \leq n/2} q^d > q^n - q^{n/2 + 1} \geq
0$（当 $n \geq 2$ 时成立，$I_1 = q \geq 1$ 亦然）：$I_n \geq 1$ 总成立。
\end{solution}

\begin{solution}{exo:b3:galois:7}
$L = \Q(\sqrt2, \sqrt3)$ 是多项式 $(X^2 -
2)(X^2 - 3)$ 的\omterm{thm:b3:galois:splitting}{分裂域}，可分，是次数为 $4$ 的\omterm{def:b3:galois:galois}{伽罗瓦扩张}
（见\cref{exo:b3:galois:1}）。自同构将 $\sqrt2 \mapsto
\pm\sqrt2$ 且 $\sqrt3 \mapsto \pm\sqrt3$：至多 $4$ 种选择，
而 $\abs G = 4$ 说明全部选择均可实现：$G \cong (\Z/2\Z)^2$，
其元素为 $\mathrm{id}, \sigma (\sqrt2 \mapsto -\sqrt2), \tau
(\sqrt3\mapsto-\sqrt3), \sigma\tau$。阶为 $2$ 的子群有：
$\langle\sigma\rangle, \langle\tau\rangle,
\langle\sigma\tau\rangle$，不动域依次为 $\Q(\sqrt3)$，
$\Q(\sqrt2)$，$\Q(\sqrt6)$（注意 $\sigma\tau$ 固定 $\sqrt6 =
\sqrt2\sqrt3$）。子域格：下方为 $\Q$，中间为三个
二次域，顶端为 $L$——且别无其他中间域
（见\cref{thm:b3:galois:fundamental}）。对于 $(X^2-2)(X^2-3)(X^2-6)$：
其\omterm{thm:b3:galois:splitting}{分裂域}仍是\emph{同一个} $L$（$\sqrt6 =
\sqrt2\sqrt3$），因此答案完全相同：\omterm{thm:b3:galois:fundamental}{伽罗瓦
对应}是\emph{扩张}的不变量，而非
为表示它而选取的多项式的不变量。
\end{solution}

\begin{solution}{exo:b3:galois:8}
(a) 由不可约性，$G$ 忠实且传递地作用于
三个根上：$G \hookrightarrow S_3$，且 $3 \mid \abs G$，故 $G
\cong A_3$ 或 $S_3$。每个 $\sigma \in G$ 置换 $x_i$，并且
$\sigma(\delta) = \varepsilon(\sigma)\,\delta$（$\delta$ 关于各根
交错的）。若 $\Delta$ 是 $\Q$ 中的平方数：
$\delta \in \Q^\times$（注意 $\delta \neq 0$，由可分性），故
$\varepsilon(\sigma) = 1$ 对所有 $\sigma$ 成立：$G \subseteq A_3$，
从而 $= A_3$。若否，则 $\delta \notin \Q$，于是存在 $\sigma$ 使
$\varepsilon(\sigma) = -1$，故 $G = S_3$。（两种情形下均有 $\Q(\delta)
= L^{G \cap A_3}$.)

(b) 在 $x_1 + x_2 + x_3 = 0$ 的条件下：$u + v = 2x_1 - (x_2 + x_3) =
3x_1$。又
\[
uv = \sum_i x_i^2 + (j + j^2)\sum_{i<k}x_ix_k
= \Bigl(\sum x_i\Bigr)^2 - 3\sum_{i<k}x_ix_k = -3p,
\]
利用 $j + j^2 = -1$ 及 $\sum_{i<k}x_ix_k = p$，则
\[
u^3 + v^3 = (u+v)^3 - 3uv(u+v) = 27x_1^3 + 9p\cdot 3x_1
= 27\,(x_1^3 + px_1) = -27q .
\]
因此 $u^3, v^3$ 是 $Y^2 + 27qY - 27p^3 = 0$ 的根
（其乘积为 $(uv)^3 = -27p^3$）：$u^3 = \frac{-27q +
\sqrt{729q^2 + 108p^3}}2$，且 $x_1 = \frac{u + v}3$，其中 $v =
-3p/u$：此即卡尔达诺公式。$S_3$ 的可解性骨架如下：
其塔为 $\Q \subseteq \Q(\delta) \subseteq \Q(\delta, j, u)$，
先添加一个平方根（$\delta$，$A_3$ 的不动域：
即 $S_3 \to S_3/A_3$ 这一步），再添加一个立方根（$u$，因 $u^3
\in \Q(\delta, j)$，即 $A_3 \to \{e\}$ 这一步）——\omterm{def:b3:groups:derived}{导出
列} $S_3 \supset A_3 \supset \{e\}$ 的具体化。
\end{solution}

\begin{solution}{exo:b3:galois:9}
$\eta_0 + \eta_1 = \zeta_5 + \zeta_5^2 + \zeta_5^3 + \zeta_5^4 =
-1$（所有 $5$ 次单位根之和为 $0$）。$\eta_0\eta_1 =
(\zeta + \zeta^4)(\zeta^2 + \zeta^3) = \zeta^3 + \zeta^4 +
\zeta^6 + \zeta^7 = \zeta^3 + \zeta^4 + \zeta + \zeta^2 = -1$
（指标按模 $5$ 计算）。于是 $\eta_0, \eta_1$ 是 $Y^2 + Y
- 1$：$\frac{-1 \pm \sqrt5}2$。由于 $\eta_0 =
2\cos\frac{2\pi}5 > 0$：$\eta_0 = \frac{\sqrt5 - 1}2$，从而
$\cos\frac{2\pi}5 = \frac{\sqrt5 - 1}4$，且 $\eta_1 =
\frac{-1-\sqrt5}2$。群
$\operatorname{Gal}(\Q(\zeta_5)/\Q) \cong (\Z/5\Z)^\times$ 是
$4$ 阶循环群：它恰有\emph{唯一}一个阶为
$2$ 的子群（$\{\pm 1\}$，即 $\zeta \mapsto \zeta^{\pm1}$），因此
$\Q(\zeta_5)$ 有唯一的二次子域
（见\cref{thm:b3:galois:fundamental}），它包含 $\eta_0
\notin \Q$，故它是 $\Q(\eta_0) = \Q(\sqrt5)$。可构造性：
$\cos\frac{2\pi}5$ 属于二次塔 $\Q \subseteq
\Q(\sqrt5)$，而 $\zeta_5$ 再向上一步二次扩张：
由\cref{thm:b3:galois:wantzel}，可作出正五边形。
\end{solution}

\begin{solution}{exo:b3:galois:10}
(a) 记 $s = S^{1/p}$，$t = T^{1/p}$（取定\omterm{def:b3:galois:closure}{代数闭包}中的
元素，满足 $s^p = S$，$t^p = T$）。$X^p - S$ 在
$K = \mathbb F_p(S, T) = (\mathbb
F_p(T))(S)$ 的分式域上不可约：对\omterm{def:b3:rings:divisibility}{素元} $S$，
在\omterm{def:b3:rings:pidufd}{唯一分解整环} $\mathbb F_p(T)[S]$ 中应用\omterm{thm:b3:rings:criteria}{艾森斯坦判别法}（见\cref{thm:b3:rings:criteria}）。于是
$[K(s):K] = p$；同理，$X^p - T$ 关于 $T$ 在
$K(s) = \mathbb F_p(s)(T)$ 的分式域上满足\omterm{thm:b3:rings:criteria}{艾森斯坦判别法}——$T$ 在
$\mathbb F_p(s)[T]$ 中仍为\omterm{def:b3:rings:divisibility}{素元}——由此得 $[L : K(s)] = p$，且 $[L:K] =
p^2$。对 $\alpha \in L$，$L = K[s, t]$，故 $\alpha = \sum
c_{ij}s^it^j$（$c_{ij} \in K$），并由\omterm{thm:b3:galois:finitefields}{弗罗贝尼乌斯
自同态}，$\alpha^p = \sum c_{ij}^p S^iT^j \in K$。

(b) 若 $L = K(\alpha)$，则 $[K(\alpha):K] = p^2$；但
$\alpha^p = a \in K$ 意味着 $\alpha$ 是 $X^p - a$ 的根，故
$\deg\pi_\alpha \leq p < p^2$，矛盾。不存在
本原元：\cref{thm:b3:galois:primitive} 确实需要
可分性（这里每个 $\pi_\alpha$ 都整除某个 $X^p - a =
(X - \alpha)^p$：纯不可分）。
\end{solution}

\begin{solution}{exo:b3:galois:11}
(a) 对 $2$ 应用\omterm{thm:b3:rings:criteria}{艾森斯坦判别法}（$2 \mid 4, 2$；$4 \nmid 2$）：$P$
不可约。若 $\alpha$ 为根，则 $[\Q(\alpha):\Q] = 5$
整除 $[L:\Q] = \abs G$（$L$ 为\omterm{thm:b3:galois:splitting}{分裂域}）。由柯西定理
（见\cref{thm:b3:groups:cauchy}）给出一个阶为 $5$ 的元素，它属于
$G \leq S_5$；在 $S_5$ 中，只有 $5$-轮换的阶为 $5$（阶
等于各轮换长度的最小公倍数）。

(b) $P'(x) = 5x^4 - 4$ 在 $\pm(4/5)^{1/4} \approx \pm
0.946$ 处为零：先为局部极大值，再为局部极小值。取值：$P(-2)
= -22 < 0$，$P(0) = 2 > 0$，$P(1) = -1 < 0$，$P(2) = 26 > 0$：
三次变号，至多有三个实根（两个临界
点）：恰有 $3$ 个实根，从而有一对
共轭复根。取\omterm{thm:b3:galois:splitting}{分裂域} $L$ 位于 $\C$ 内：
复共轭把 $L$ 映到自身（它置换根，
这些根生成 $L$），并固定 $\Q$，故它确定了
$G$ 中的一个元素：它固定三个实根，交换另外两根，是一个
对换。

(c) 设 $\tau = (a\,b)$，$\sigma$ 为 $5$-轮换，且属于 $G$。某个
幂 $\sigma^k$ 把 $a$ 映到 $b$（$k \ne 0 \bmod 5$），而
$\sigma^k$ 仍是 $5$-轮换。适当改记号，可设 $\sigma =
(1\,2\,3\,4\,5)$，$\tau = (1\,2)$。作共轭，
$\sigma^m\tau\sigma^{-m} = (\sigma^m(1)\ \sigma^m(2))$：
相邻对换 $(1\,2), (2\,3), (3\,4), (4\,5),
(5\,1)$ 均属于 $G$。相邻对换生成 $S_5$
（每个对换 $(i\,j)$ 皆为相邻对换之积，
而对换生成全体置换）。因此 $G = S_5$，不可解
（见\cref{cor:b3:groups:snnotsolvable}），再由
\cref{thm:b3:galois:solvable} 得出：$X^5 - 4X + 2$
根式不可解。
\end{solution}

\begin{solution}{exo:b3:galois:12}
(a) $X^4 - 2$ 不可约（由 $2$ 处的\omterm{thm:b3:rings:criteria}{艾森斯坦判别法}）：
$[\Q(\alpha):\Q] = 4$；$\iu \notin \Q(\alpha) \subseteq \R$，
故 $[L : \Q(\alpha)] = 2$，$[L:\Q] = 8$。该扩张
是\omterm{def:b3:galois:galois}{伽罗瓦扩张}（\omterm{def:b3:galois:separable}{可分多项式}的\omterm{thm:b3:galois:splitting}{分裂域}：根为
$\iu^k\alpha$），故 $\abs G = 8$。自同构把
$\alpha$ 映到四个根之一，$\iu$ 映到 $\pm\iu$：
至多 $8$ 个映射，且均可实现。给定的 $\sigma$（阶
为 $4$：$\sigma^2(\alpha) = -\alpha$，$\sigma^4 = e$）与 $\tau$
（阶为 $2$）满足
\[
\tau\sigma\tau(\alpha) = \tau\sigma(\alpha) =
\tau(\iu\alpha) = -\iu\alpha = \sigma^{-1}(\alpha),
\qquad \tau\sigma\tau(\iu) = \iu\cdot(-1)(-1) = \iu ,
\]
更细致地：$\tau\sigma\tau(\iu) = \tau\sigma(-\iu) =
\tau(-\iu) = \iu = \sigma^{-1}(\iu)$。于是 $\tau\sigma\tau =
\sigma^{-1}$，即 $D_4$ 的群表现。

(b) $D_4 = \langle\sigma,
\tau\rangle$ 的十个子群：$\{e\}$；五个阶为 $2$ 的
$\langle\sigma^2\rangle$，$\langle\tau\rangle$，
$\langle\sigma^2\tau\rangle$，$\langle\sigma\tau\rangle$，
$\langle\sigma^3\tau\rangle$；三个阶为 $4$ 的：
$\langle\sigma\rangle$，$\{e, \sigma^2, \tau,
\sigma^2\tau\}$，$\{e, \sigma^2, \sigma\tau,
\sigma^3\tau\}$；以及 $D_4$。不动域（次数等于指数）：
$\{e\} \leftrightarrow L$；阶为 $2$ 的子群
$\leftrightarrow$ 五个四次域
\[
\langle\tau\rangle \leftrightarrow \Q(\alpha),\quad
\langle\sigma^2\tau\rangle \leftrightarrow \Q(\iu\alpha),
\quad
\langle\sigma^2\rangle \leftrightarrow \Q(\sqrt2, \iu),\quad
\langle\sigma\tau\rangle \leftrightarrow
\Q\bigl((1+\iu)\alpha\bigr),\quad
\langle\sigma^3\tau\rangle \leftrightarrow
\Q\bigl((1-\iu)\alpha\bigr) .
\]
验证：$\tau$ 固定实数 $\alpha$；$\sigma^2\tau$ 把
$\alpha \mapsto -\alpha$，$\iu \mapsto -\iu$，并固定
$\iu\alpha$；又由 $\sigma\tau(\alpha) = \iu\alpha$，
$\sigma\tau(\iu) = -\iu$：
\[
\sigma\tau\bigl((1+\iu)\alpha\bigr) =
(1 - \iu)\,\iu\alpha = (1 + \iu)\alpha,
\qquad
\sigma^3\tau\bigl((1-\iu)\alpha\bigr) =
(1 + \iu)(-\iu)\alpha = (1 - \iu)\alpha :
\]
每个反射都固定其生成元，且其不动域
（次数 $4 =$ 指数）恰为它所生成的域
（生成元为 $X^4 + 8$ 的根，不可约）。阶为 $4$ 的子群
$\leftrightarrow$ 三个二次域：
$\langle\sigma\rangle \leftrightarrow \Q(\iu)$（$\sigma$
固定 $\iu$）；$\{e, \sigma^2, \tau, \sigma^2\tau\}
\leftrightarrow \Q(\sqrt2)$（四者都固定 $\alpha^2$，至多相差
符号；验证：$\tau(\sqrt2) = \sqrt2$，$\sigma^2(\alpha^2) =
(-\alpha)^2$）；$\{e, \sigma^2, \sigma\tau, \sigma^3\tau\}
\leftrightarrow \Q(\iu\sqrt2)$（$\sigma\tau(\iu\alpha^2) =
(-\iu)(\iu\alpha)^2 = \iu\alpha^2$）。

(c) 在 $\Q$ 上伽罗瓦的中间域 $\leftrightarrow$
$D_4$ 的\omterm{def:b3:groups:normal}{正规子群}：$\{e\}$，$\langle\sigma^2\rangle$（中心），
阶为 $4$ 的三个子群，以及 $D_4$——因此伽罗瓦中间域为
$L$，$\Q(\sqrt2, \iu)$，三个
二次域，以及 $\Q$。由非正规反射所固定的五个
四次域均不是\omterm{def:b3:galois:galois}{伽罗瓦扩张}：$\Q(\alpha)$ 含有
$X^4 - 2$ 的一个根，但不含 $\iu\alpha$（它是实数）——
$\sigma$ 的共轭把 $\langle\tau\rangle$ 变为
$\langle\sigma^2\tau\rangle$，正如它把
$\Q(\alpha)$ 变为 $\Q(\iu\alpha)$：子群的非正规性
\emph{正是}存在一个共轭域。
\end{solution}

\begin{solution}{pb:b3:galois:1}
\textbf{1.} $\Phi_{17}$ 不可约
（见\cref{thm:b3:galois:cyclotomicirred}，或
\cref{ex:b3:rings:eisenstein}，当指数为素数时）：$[L:\Q] =
\varphi(17) = 16$，且 $G \cong (\Z/17\Z)^\times$，是
$16$ 阶循环群（见\cref{thm:b3:galois:cyclic}）。$3$ 的幂模
$17$：
\[
3,\ 9,\ 10,\ 13,\ 5,\ 15,\ 11,\ 16,\ 14,\ 8,\ 7,\ 4,\ 12,\ 2,\
6,\ 1
\]
——十六个互不相同的值，$3$ 为生成元。

\textbf{2.} $G = \langle\sigma\rangle$ 为 $16$ 阶循环群；
$H_k = \langle\sigma^{2^k}\rangle$ 的阶为 $2^{4-k}$，且
$[H_k : H_{k+1}] = 2$。由基本定理
（\cref{thm:b3:galois:fundamental}），$L_k = L^{H_k}$ 满足
$[L_k : \Q] = [G : H_k] = 2^k$，故每个 $[L_{k+1}:L_k] = 2$。

\textbf{3.} $\zeta \in L = L_4$ 位于 $\Q$ 的一串二次
扩张之上：由\cref{thm:b3:galois:wantzel}，$\zeta$ 可
构造；正 $17$ 边形的顶点为 $\zeta^k$。

\textbf{4.} $\sigma^2$ 将指数乘以 $9$；
$\eta_0$ 的指数为 $3$ 的偶次幂，
\[
\{3^{2k} \bmod 17\} = \{1, 9, 13, 15, 16, 8, 4, 2\},
\]
该集合在乘以 $9 = 3^2$ 下不变；故 $\eta_0$
（$\eta_1$ 同理）被 $H_1 =
\langle\sigma^2\rangle$ 固定：$\eta_0, \eta_1 \in L_1$，
为二次域。$\sigma$ 把偶次幂映为奇次幂，它交换 $\eta_0,
\eta_1$。于是 $\eta_0 + \eta_1$ 与 $\eta_0\eta_1$ 被
$G$ 中全部元素固定，故为有理数；$\eta_0, \eta_1$ 是一个
有理系数二次多项式的根。

\textbf{5.} $\eta_0 + \eta_1 = \sum_{c=1}^{16}\zeta^c = -1$。乘积
展开为 $64$ 项 $\zeta^{a + b}$，$a$ 属于偶数集，
$b$ 属于奇数集。没有哪一项是 $\zeta^0$：$b = -a$
不可能，因为 $-1 = 16 = 3^8$ 是\emph{偶}次幂，故
$-a$ 仍属于偶数集。于是 $\eta_0\eta_1 = \sum_{c \neq 0}
n_c\zeta^c$，$\sum n_c = 64$；作用 $\sigma$ 固定
$\eta_0\eta_1$（它交换两个因子），并置换 $\zeta^c$，
在全部 $c \neq 0$ 上传递地，故所有 $n_c$ 相等：$n_c
= 4$，且 $\eta_0\eta_1 = 4\sum_{c\neq0}\zeta^c = -4$。

\textbf{6.} $\eta_{0,1}$ 满足 $Y^2 + Y - 4 = 0$：$\frac{-1 \pm
\sqrt{17}}2$。数值上，配对共轭指数，
$\eta_0 = 2\bigl(\cos\tfrac{2\pi}{17} + \cos\tfrac{4\pi}{17} +
\cos\tfrac{8\pi}{17} + \cos\tfrac{16\pi}{17}\bigr) \approx 1.56
> 0$: $\eta_0 = \frac{-1+\sqrt{17}}2$, $\eta_1 =
\frac{-1-\sqrt{17}}2$，且 $L_1 = \Q(\eta_0) = \Q(\sqrt{17})$。

\textbf{7.} 指数集：$\beta_0$ 为 $\{1, 13, 16, 4\}$ =
$3^{4k}$ 的幂；$\beta_2$ 为 $\{9, 15, 8, 2\}$ = $9 \times$ 该
集合。并集即偶数集：$\beta_0 + \beta_2 = \eta_0$；同理
$\beta_1 + \beta_3 = \eta_1$。乘以 $13 = 3^4$
使每个 $\beta_i$ 的指数集不变：被 $H_2 =
\langle\sigma^4\rangle$ 固定；而 $\sigma^2$（$\times 9$）将
$\{1,13,16,4\}$ 映为 $\{9, 15, 8, 2\}$，交换 $\beta_0,
\beta_2$。

\textbf{8.} 展开 $\beta_0\beta_2$，十六个指数
之和
\[
\{1,13,16,4\} + \{9,15,8,2\} =
\{10,16,9,3,\ 5,11,4,15,\ 8,14,7,1,\ 13,2,12,6\}
\]
恰好覆盖 $1, \dots, 16$ 各一次：$\beta_0\beta_2 =
\sum_{c\ne0}\zeta^c = -1$。作用 $\sigma$（它把
$\beta_0 \mapsto \beta_1$，$\beta_2 \mapsto \beta_3$：
指数 $\times 3$），$\beta_1\beta_3 = \sigma(\beta_0\beta_2)
= -1$。

\textbf{9.} $\beta_0, \beta_2$ 满足 $Y^2 - \eta_0 Y - 1 = 0$，
故 $\beta_0 = \frac{\eta_0 + \sqrt{\eta_0^2 + 4}}2$（数值上
$\beta_0 = 2\cos\frac{2\pi}{17} + 2\cos\frac{8\pi}{17} \approx
2.05 > 0$，取 $+$ 号）。同理 $\beta_1 = \frac{\eta_1 +
\sqrt{\eta_1^2 + 4}}2 \approx 0.344$（数值检验再次确定
符号）。$L_2 = L_1(\beta_0)$，是 $L_1$ 上的二次扩张。

\textbf{10.} $\gamma_0 + \gamma_1 = \zeta + \zeta^{16} +
\zeta^{13} + \zeta^4 = \beta_0$。又有
\[
\gamma_0\gamma_1 = (\zeta + \zeta^{16})(\zeta^{13} + \zeta^4)
= \zeta^{14} + \zeta^{5} + \zeta^{12} + \zeta^{3} = \beta_1 .
\]
故 $\gamma_0, \gamma_1$ 满足 $Y^2 - \beta_0Y + \beta_1 = 0$；
数值上 $\gamma_0 = 2\cos\frac{2\pi}{17} \approx 1.865 >
\gamma_1 \approx 0.185$：$\gamma_0 = \frac{\beta_0 +
\sqrt{\beta_0^2 - 4\beta_1}}2$。

\textbf{11.} 逐步代入：
\[
\eta_0 = \frac{-1 + \sqrt{17}}2,
\quad
\beta_0 = \frac{\eta_0 + \sqrt{\eta_0^2 + 4}}2,
\quad
\beta_1 = \frac{\eta_1 + \sqrt{\eta_1^2 + 4}}2,
\quad
\cos\frac{2\pi}{17} = \frac{\beta_0 + \sqrt{\beta_0^2 -
4\beta_1}}4 .
\]
数值上：$\sqrt{17} \approx 4.1231$，$\eta_0 \approx
1.5616$，$\eta_1 \approx -2.5616$，$\beta_0 \approx 2.0494$，
$\beta_1 \approx 0.3441$，$\beta_0^2 - 4\beta_1 \approx
2.8234$，且 $\cos\frac{2\pi}{17} \approx \frac{2.0494 +
1.6803}4 \approx 0.93242$——与 $\cos\frac{2\pi}{17} =
0.93247\dots$：微小差异源于
中间式子中的舍入；多取几位即可复现
$0.932472$。

\textbf{12.} 该构造要求 $[L:\Q] = p - 1$ 为
$2$ 的幂，使得由指数为 $2$ 的子群所成的完整链
存在。若 $p = 2^m + 1$ 是质数，且 $m = ab$，其中 $a$ 为奇数且 $>
1$：$x + 1 \mid x^a + 1$ 在 $x = 2^b$ 处表明 $2^b + 1$ 真
整除 $p$——不可能。于是 $m$ 为 $2$ 的幂：$p =
2^{2^t} + 1$，即\emph{费马质数}（$3, 5, 17, 257, 65537$，
\dots）。反之，对此类 $p$，$\varphi(p) = 2^{2^t}$，且
沿用 1--3 题的论证（或
\cref{cor:b3:galois:impossible}）即得：正 $p$ 边形
可构造，当且仅当 $p$ 为费马质数。

\textbf{13.} $\varphi(n)$ 为 $2$ 的幂，当且仅当 $n =
2^a p_1\cdots p_r$，其中 $p_i$ 为互异的费马质数
（$\varphi$ 的乘性；奇素数幂 $p^k$，$k
\geq 2$，贡献因子 $p \nmid 2^m$）。对 $n \leq
20$，\omterm{def:b3:galois:constructible}{可构造的}正 $n$ 边形为
\[
n = 3, 4, 5, 6, 8, 10, 12, 15, 16, 17, 20
\]
相应的取值为
\[
\varphi(n) = 2,\ 2,\ 4,\ 2,\ 4,\ 4,\ 4,\ 8,\ 8,\ 16,\ 8 .
\]
不\omterm{def:b3:galois:constructible}{可构造的}情形为 $n = 7, 9, 11, 13, 14, 18, 19$，其中
$\varphi(n) = 6, 6, 10, 12, 6, 6, 18$ 含有奇素因子。

\textbf{14.} 平方剩余构成平方映射的像，
该映射定义于循环群 $(\Z/p\Z)^\times$ 上，所得像的指数为 $2$：
共 $\frac{p-1}2$ 个平方剩余，$\frac{p-1}2$ 个非平方剩余，故
勒让德符号之和为 $0$。于是
\[
\sum_{a=0}^{p-1}\zeta^{a^2} = 1 + 2\sum_{b \text{ 平方}
\neq 0}\zeta^b
= 1 + \sum_{b\neq0}\Bigl(1 + \Bigl(\frac bp\Bigr)\Bigr)
\zeta^b = \sum_{b}\zeta^b + g = g,
\]
利用 $1 + \bigl(\frac bp\bigr) = \#\{a : a^2 = b\}$ 及
$\sum_{b=0}^{p-1}\zeta^b = 0$。对 $p = 17$，模
$17$ 的平方剩余即生成元 $3$ 的偶次幂，亦即
$\eta_0$（第二部分）中出现的指数，故 $g =
\sum_{\text{偶数 }k}\zeta^{3^k} - \sum_{\text{奇数 }k}
\zeta^{3^k} = \eta_0 - \eta_1$。

\textbf{15.} 令 $b = c - a$（$a, b$ 取遍非零
剩余，$c = a + b$ 取遍全体剩余）：
\[
g^2 = \sum_c\zeta^c\sum_{a\neq0,\,a\neq c}
\Bigl(\frac{a(c-a)}p\Bigr) .
\]
当 $c = 0$ 时，$\bigl(\frac{-a^2}p\bigr) =
\bigl(\frac{-1}p\bigr)$，对 $p - 1$ 个值求和。当 $c \neq 0$ 时，作代换 $c - a = at$，即 $t = c/a -
1$；当 $a$ 取遍非零剩余时，$t$
双射地取遍 $\neq -1$ 的全体剩余（反解：$a = c/(1 +
t)$）。求和项化为 $\bigl(\frac{a^2t}p\bigr) =
\bigl(\frac tp\bigr)$，且
\[
\sum_{t \neq -1}\Bigl(\frac tp\Bigr)
= -\Bigl(\frac{-1}p\Bigr)
\]
（由第 14 题，全和为零）。因此
\[
g^2 = \Bigl(\frac{-1}p\Bigr)\Bigl[(p-1) -
\sum_{c\neq0}\zeta^c\Bigr]
= \Bigl(\frac{-1}p\Bigr)\,p = p^* ,
\]
利用 $\sum_{c\neq0}\zeta^c = -1$ 及欧拉判别法
$\bigl(\frac{-1}p\bigr) = (-1)^{(p-1)/2}$。当 $p = 17$ 时：
$(\eta_0 - \eta_1)^2 = (\eta_0 + \eta_1)^2 - 4\eta_0\eta_1 =
1 + 16 = 17$，与第二部分相符。

\textbf{16.} $g^2 = p^*$ 表明 $\sqrt{p^*} = \pm g \in
\Q(\zeta_p)$，故 $\Q(\sqrt{p^*})$ 为二次子域。
唯一性：次数为 $2$ 的子域，由\omterm{thm:b3:galois:fundamental}{伽罗瓦对应}，
恰与指数为 $2$ 的子群一一对应，这些子群位于
循环群 $\operatorname{Gal}(\Q(\zeta_p)/\Q) \cong (\Z/p\Z)^\times$ 中，
而偶数阶循环群恰有一个此类子群
（即平方剩余所成的子群）。每个二次域都形如
$\Q(\sqrt{d})$，其中 $d$ 无平方因子。将域
$\Q(\sqrt{p^*})$、$\Q(\iu) \subseteq \Q(\zeta_4)$ 及
$\Q(\sqrt2) \subseteq \Q(\zeta_8)$ 同含于某一公共的
$\Q(\zeta_N)$ 中，即可涵盖所有 $\sqrt d$，
此即克罗内克--韦伯定理的二次情形。

\textbf{17.} 在任一交换环中，$(x + y)^q = x^q + y^q
+ q(\cdots)$：二项式系数 $\binom qk$，$0 < k <
q$，均可被素数 $q$ 整除。对 $p - 1$
项之和 $g$ 反复应用上式，得
\[
g^q \equiv \sum_{a}\Bigl(\frac ap\Bigr)^{q}\zeta^{aq}
= \sum_a\Bigl(\frac ap\Bigr)\zeta^{aq}
\pmod{q\Z[\zeta]},
\]
（$q$ 为奇数时，符号不变）。作替换 $b = aq$，$a =
q^{-1}b$，且 $\bigl(\frac{q^{-1}b}p\bigr) =
\bigl(\frac{q}p\bigr)\bigl(\frac bp\bigr)$（乘性；
$\bigl(\frac{q^{-1}}p\bigr) = \bigl(\frac qp\bigr)$，因逆元的
符号与原元相同）：$g^q \equiv
\bigl(\frac qp\bigr)g$。

\textbf{18.} $g^q = g\,(g^2)^{(q-1)/2} = g\,(p^*)^{(q-1)/2}$
精确成立（由第 15 题），并由 $\Z$ 中的欧拉判别法给出
$(p^*)^{(q-1)/2} \equiv \bigl(\frac{p^*}q\bigr) \pmod q$,
故模 $q\Z[\zeta]$ 有 $g^q \equiv \bigl(\frac{p^*}q\bigr)g$。
与第 17 题比较并乘以 $g$：
\[
\Bigl(\frac qp\Bigr)p^* \equiv \Bigl(\frac{p^*}q\Bigr)p^*
\pmod{q\Z[\zeta]} .
\]
两边均为有理整数；其差为 $0$ 或
$\pm2p^*$，属于 $q\Z[\zeta] \cap \Z = q\Z$（整数 $m
\in q\Z[\zeta]$ 满足 $m/q \in \Q \cap \Z[\zeta] = \Z$，
其中因 $1, \zeta, \dots, \zeta^{p-2}$ 构成
$\Q$ 上的基，其有理坐标可读出
整性）。由于 $q \nmid 2p^*$（$q$ 为奇数，$q \neq p$），
故其差为 $0$：$\bigl(\frac qp\bigr) =
\bigl(\frac{p^*}q\bigr)$。

\textbf{19.} 由乘性，$\bigl(\frac{p^*}q\bigr) =
\bigl(\frac{-1}q\bigr)^{(p-1)/2}\bigl(\frac pq\bigr) =
(-1)^{\frac{q-1}2\cdot\frac{p-1}2}\bigl(\frac pq\bigr)$,
故第 18 题化为 $\bigl(\frac qp\bigr)\bigl(\frac
pq\bigr) = (-1)^{\frac{p-1}2\frac{q-1}2}$：即互反律。
验证 $(17, 3)$：指数 $\frac{16}2\cdot\frac22 = 8$
为偶数，故两符号须相等；模 $3$ 的平方剩余为 $\{1\}$，而 $17 \equiv 2$：
$\bigl(\frac{17}3\bigr) = -1$；模 $17$ 的平方剩余为 $\{1, 4,
9, 16, 8, 2, 15, 13\}$，而 $3$ 不在其中：
$\bigl(\frac3{17}\bigr) = -1$。乘积为 $+1$，与预期一致。对 $x^2
\equiv 219 \pmod{383}$: $\bigl(\frac{219}{383}\bigr) =
\bigl(\frac3{383}\bigr)\bigl(\frac{73}{383}\bigr)$。首先
$383 \equiv 3 \pmod4$ 且 $3 \equiv 3$：由互反律得
$\bigl(\frac3{383}\bigr) = -\bigl(\frac{383}3\bigr) =
-\bigl(\frac23\bigr) = -(-1) = +1$。其次，$73 \equiv 1
\pmod 4$: $\bigl(\frac{73}{383}\bigr) =
\bigl(\frac{383}{73}\bigr) = \bigl(\frac{18}{73}\bigr) =
\bigl(\frac2{73}\bigr)$（$18 = 2\cdot3^2$），且 $73 \equiv 1
\pmod 8$ 使 $2$ 为模 $73$ 的平方剩余（补充律，
可仿上法证明，取 $g = \zeta_8 + \zeta_8^{-1} = \sqrt2$，在
$\Q(\zeta_8)$ 中同法证明）：$+1$。总计 $+1$，
同余式可解。

\textbf{20.} 本原根的存在性与唯一性：
若 $w$ 的周期集合为 $\{d : w = u^{n/d},\ \abs u = d\}$，
则满足该条件的最小者 $d_0$ 整除其他每个周期 $d$
（若 $w$ 既是 $d$ 次幂又是 $d'$ 次幂，则它是
$\gcd(d, d')$ 次幂：比较模其最大公因数同余的下标处的字母，
借助贝祖定理），且长为 $d_0$ 的块是
本原的。将 $q^n$ 个字按其本原根的长度
分类：$q^n = \sum_{d \mid n}A(d)$。由于 $A$ 与
$d\,N_q(d)$ 满足相同的递推关系，并
在 $n = 1$ 时取值相同（二者均可由 $q^n
= \sum_{d\mid n}(\cdot)$ 归纳互定），因此二者相等：$A(d) =
d\,N_q(d)$。双射：元素 $\alpha$ 的次数为 $d$，
对应字 $w_\alpha$，即其系数的\dots\
不如直接：次数为 $d$ 的元素
位于 $\overline{\mathbb F_q}$ 中者，恰为 $N_q(d)$
个次数为 $d$ 的不可约多项式的根，每个贡献其 $d$
个互异的根（可分性）：共 $d\,N_q(d)$ 个
次数为 $d$ 的元素，与计数 $q^n = \sum_{d\mid n}\#\{
\text{ 次数为 } d \text{ 的元素，位于 } \mathbb F_{q^n}\}$ 相符，
——同一筛法，一次作用于字，一次作用于域元素。

\textbf{21.} 引理：$\sum_{d \mid m}\mu(d) = \mathbf
1_{m=1}$。对 $m > 1$，固定素数 $p \mid m$：
$m$ 的无平方因数可两两配对为 $\{d, pd\}$，其中 $p \nmid d$，
且 $\mu(pd) = -\mu(d)$，故和为零。于是，对 $f(n) =
\sum_{e\mid n}g(e)$：
\[
\sum_{d \mid n}\mu(d)\,f\bigl(\tfrac nd\bigr)
= \sum_{d \mid n}\mu(d)\!\!\sum_{e \mid n/d}\!\!g(e)
= \sum_{e \mid n}g(e)\!\!\sum_{d \mid n/e}\!\!\mu(d)
= g(n) .
\]
取 $f(n) = q^n$ 及 $g(n) = nN_q(n)$
（见\cref{exo:b3:galois:6}），得 $N_q(n) =
\frac1n\sum_{d\mid n}\mu(d)\,q^{n/d}$.

\textbf{22.} $d = 1$ 的项为 $q^n$；其余每一项都
满足 $\abs{\mu(d)q^{n/d}} \leq q^{n/2}$，且粗略地
$\sum_{d \mid n, d > 1}q^{n/d} \leq \sum_{j \leq
n/2}q^j < 2q^{n/2}$（几何级数，$q \geq 2$）。于是 $nN_q(n) >
q^n - 2q^{n/2} \geq 0$ 对 $n \geq 1$ 成立：每个次数的不可约多项式
都存在，且 $\mathbb F_q[X]/(\pi) = \mathbb F_{q^n}$ 可
（重新）构造——以计数方式证明存在性。
在 $q^n$ 个首一 $n$ 次多项式中，不可约者所占比例为
$\frac1n(1 + O(q^{-n/2}))$：这即
$\mathbb F_q[X]$ 上的素数定理，其中 $n$ 扮演 $\log x$ 的角色。当 $q = 2$ 时，
\cref{exo:b3:galois:6} 中的计数 $2, 1, 2, 3$ 与公式相符：
例如 $N_2(4) = \frac14(2^4 - 2^2) = 3$。

\textbf{23.} 在 $\mathbb F_q$ 上非零有理函数构成的乘法交换群中，
令 $F(n) = X^{q^n} -
X$ 及 $G(n) = \prod_{\deg\pi = n}\pi$；
由 \cref{exo:b3:galois:6} 得 $F(n) = \prod_{d\mid n}G(d)$。
第 21 题中的莫比乌斯论证，
写成乘法形式（指数的加法恰与原来的求和一致），
得 $G(n) = \prod_{d \mid n}F(d)^{\mu(n/d)}$。当 $q = 2$，
$n = 2$ 时，$G(2) = \frac{X^4 - X}{X^2 - X} = \frac{X(X^3 -
1)}{X(X - 1)} = X^2 + X + 1$，即
$\mathbb F_2$ 上唯一的不可约二次多项式，符合预期。

\textbf{24.} $\omega^2 = \iu$ 且 $\omega^{-2} = -\iu$，故
$g^2 = \omega^2 + 2 + \omega^{-2} = 2$。在交换环
$\Z[\omega]/q\Z[\omega]$ 中应用“新生梦”：$g^q =
(\omega + \omega^{-1})^q \equiv \omega^q + \omega^{-q}$。
$\omega^q + \omega^{-q}$ 的值仅取决于 $q \bmod 8$：
当 $q \equiv \pm1$ 时，$\omega^q + \omega^{-q} = \omega^{\pm1} +
\omega^{\mp1} = g$；当 $q \equiv \pm3$ 时，利用 $\omega^4 =
-1$，$\omega^{3} = -\omega^{-1}$ 及 $\omega^{-3} = -\omega$，
故 $\omega^q + \omega^{-q} = -g$。另一方面，$g^q =
g\,(g^2)^{(q-1)/2} = g\,2^{(q-1)/2} \equiv \bigl(\frac2q\bigr)
g \pmod{q\Z[\omega]}$，由模 $q$ 的欧拉判别法。比较
两边并乘以 $g$：$2\bigl(\frac2q\bigr) \equiv
\pm2 \pmod{q\Z[\omega]}$；若两者符号不同，则 $q$ 将
整除 $4$，在 $\Z[\omega]$ 中，从而在 $\Z$ 中亦然（$q\Z[\omega] \cap
\Z = q\Z$，按基 $1, \omega, \omega^2,
\omega^3$ 的坐标），这对奇数 $q$ 不可能。于是 $\bigl(\frac2q\bigr) =
+1$ 当且仅当 $q \equiv \pm1 \pmod 8$。奇偶性检验：$q = 8k \pm 1$
给出 $(q^2 - 1)/8 = 2k(4k \pm 1)$，为偶数；$q = 8k \pm 3$
给出 $(q^2 - 1)/8 = 8k^2 \pm 6k + 1$，为奇数：公式
$(-1)^{(q^2-1)/8}$ 正好刻画了这一分情形。数值上：$3^2
= 9 \equiv 2 \pmod 7$（$7 \equiv -1$），$6^2 = 36 \equiv 2
\pmod{17}$（$17 \equiv 1$）；模 $3$ 的平方数为 $\{0,
1\}$，模 $5$ 的为 $\{0, 1, 4\}$，二者均不含 $2$
（$3 \equiv 3$，$5 \equiv -3 \pmod 8$）。

\textbf{25.} 每个首一的 $f \in \mathbb F_q[X]$ 唯一地
写成首一不可约多项式的幂之积 $\prod_\pi\pi^{e_\pi}$，
按次数归类，
\[
\sum_{f \text{ 首一}}t^{\deg f}
= \prod_\pi\ \sum_{e \geq 0}t^{e\deg\pi}
= \prod_\pi\bigl(1 - t^{\deg\pi}\bigr)^{-1}
= \prod_{n \geq 1}\bigl(1 - t^n\bigr)^{-N_q(n)},
\]
所有乘积在 t-进意义下都合理（只有次数 $\leq m$
的项影响 $t^m$ 的系数，且每个次数的不可约多项式
只有有限个）。左边为 $\sum_m
q^mt^m = (1 - qt)^{-1}$：恒等式。取对数：$-\log(1 -
qt) = \sum_m\frac{q^m}mt^m$，而 $\sum_nN_q(n)\bigl(-\log(1
- t^n)\bigr) = \sum_nN_q(n)\sum_k\frac{t^{nk}}k$；比较
$t^m$ 的系数，得 $\frac{q^m}m = \sum_{dk = m}
\frac{N_q(d)}k = \frac1m\sum_{d \mid m}d\,N_q(d)$，即
$q^m = \sum_{d\mid m}d\,N_q(d)$。手工验算，$q = 2$，
$t^2$ 的系数：$N_2(1) = 2$，$N_2(2) = 1$，且
$(1 - t)^{-2}(1 - t^2)^{-1} = (1 + 2t + 3t^2 + \dots)(1 + t^2
+ \dots)$ 的 $t^2$ 系数为 $3 + 1 = 4 = 2^2$。最后，
利用第 21 题的公式，取因数 $1, 2, 3, 6$：
\[
N_2(6) = \tfrac16\bigl(2^6 - 2^3 - 2^2 + 2\bigr)
= \tfrac{54}6 = 9,
\]
确实有 $1\cdot2 + 2\cdot1 + 3\cdot2 + 6\cdot9 = 2 + 2 +
6 + 54 = 64 = 2^6$。
\end{solution}
